\documentclass[11pt]{article}
\usepackage{amsmath}
\usepackage{amssymb}
\usepackage{booktabs}
\usepackage{graphicx}
\usepackage{hyperref}
\usepackage{natbib}
\usepackage{geometry}
\geometry{a4paper, margin=1in}

\title{Dynamic Optimal Leverage Under Volatility Regimes:\\A Regime-Switching Framework for Risk Allocation}
\author{Quantitative Research Team}
\date{\today}

\begin{document}

\maketitle

\begin{abstract}
This paper presents a regime-switching framework to dynamically allocate leverage in equity portfolios. Using the historical S\&P 500 index dataset from Robert Shiller spanning over a century, we study the transition patterns of volatility and investigate whether dynamic adjustment of leverage can enhance risk-adjusted performance compared to constant-leverage strategies and traditional volatility-targeting strategies. We employ a Gaussian Mixture Model (GMM) with five distinct regimes representing Low Volatility, Normal Market, High Volatility, Crisis, and Recovery. We demonstrate that regimes are highly persistent. Furthermore, we develop a Recovery Acceleration leverage model that incorporates momentum, volatility decline, and drawdown recovery speed to safely enhance portfolio compounding during recovery periods. By checking robustness via bootstrap simulations, we show that regime-dependent dynamic leverage strategies statistically outperform traditional strategies on a Sharpe and Sortino ratio basis, while dramatically reducing peak-to-trough drawdowns during severe historical crises.
\end{abstract}

\section{Introduction}
Leverage is a double-edged sword in quantitative finance. While it scales positive returns, it scaling losses linearly and impairs long-term compounding through volatility drag. Traditional fixed leverage (e.g., constant 2x exposure) suffers from extreme drawdowns during market crises, often leading to margin calls or total capital loss. Volatility targeting is a common alternative that scales leverage inversely with rolling realized volatility:
\begin{equation}
L_t = \frac{\sigma_{\text{target}}}{\hat{\sigma}_t}
\end{equation}
However, simple volatility targeting assumes a continuous relationship between volatility and returns, neglecting the distinct structural changes that occur when markets transition into crises or recovery phases.

In this paper, we propose a regime-switching framework that models market states explicitly. We raise the core research question: \textit{Can a volatility-regime-based dynamic leverage strategy improve long-term risk-adjusted returns while reducing drawdowns compared with traditional fixed leverage and volatility targeting?}

Our contribution is twofold. First, we identify five distinct regimes based on multi-period volatility and drawdown features using a Gaussian Mixture Model. Second, we design and backtest a recovery-aware leverage model that gradually accelerates exposure during recovery phases to exploit the return drift typical of post-crisis expansions, while drastically scaling down exposure during high-volatility and crisis states.

\section{Literature Review}
The modeling of volatility regimes has a rich history in empirical finance. Hamilton's seminal work on Markov-switching models introduced the idea that economic variables undergo discrete state transitions \citep{hamilton1989}. Volatility regimes are well-documented to correspond to different risk premium regimes \citep{ang2002}. In periods of low volatility, returns are typically steady and positive, whereas high-volatility states are frequently accompanied by large negative returns \citep{french1987}.

Volatility targeting has been popularized in recent years as a method to manage portfolio risk. Moreira and Muir show that volatility-managed portfolios produce substantial utility gains for investors by timing volatility \citep{moreira2017}. However, their focus is on linear volatility scaling rather than discrete regime-based scaling. 

Kritzman et al. highlight that regime-switching frameworks can improve asset allocation, as the correlation structure and covariance matrix of assets change during crisis periods \citep{kritzman2012}. During crisis regimes, there is often a ``flight to safety'' where risky asset returns become highly correlated and drop precipitously \citep{baele2019}. The recovery phase after a crisis represents a unique regime characterized by declining volatility and strong positive returns, which can be leveraged to accelerate recovery speed.

\section{Data}
We use the monthly S\&P 500 dataset compiled by Robert Shiller, spanning from January 1871 onwards \citep{shiller2015}. The dataset includes the S\&P price index ($P$), dividends ($D$), earnings ($E$), and the Consumer Price Index ($CPI$).

The primary variable of interest is the price series $P$. We calculate monthly log returns as:
\begin{equation}
r_t = \log\left(\frac{P_t}{P_{t-1}}\right)
\end{equation}
The dataset is clean and historical, offering a unique opportunity to evaluate our strategies across multiple market cycles, including the Great Depression (1929), the post-WWII reconstruction, the 1970s stagflation, the 1987 crash, the 2000 Dot-com bust, the 2008 Global Financial Crisis, and the 2020 COVID-19 pandemic crash.

\section{Methodology}

\subsection{Volatility Regime Model}
We model market regimes using a Gaussian Mixture Model (GMM) with 5 clusters. Rather than relying on simple price indicators, we construct a feature vector $\mathbf{x}_t$ for each month $t$ containing:
\begin{enumerate}
    \item 12-month rolling realized volatility: $\sigma_{12, t}$
    \item 3-month change in realized volatility: $\Delta\sigma_{12, t} = \sigma_{12, t} - \sigma_{12, t-3}$
    \item Drawdown from the expanding historical peak: $DD_t = \frac{P_t}{\max_{s \le t} P_s} - 1$
\end{enumerate}

The GMM models the probability density of the features as a mixture of multivariate Gaussians:
\begin{equation}
p(\mathbf{x}_t) = \sum_{k=1}^{K} \pi_k \mathcal{N}(\mathbf{x}_t; \boldsymbol{\mu}_k, \boldsymbol{\Sigma}_k)
\end{equation}
where $\pi_k$ represents the mixing weight of regime $k$, and $\boldsymbol{\mu}_k, \boldsymbol{\Sigma}_k$ represent the mean vector and covariance matrix respectively. We fit the GMM using the Expectation-Maximization (EM) algorithm \citep{dempster1977}.

The five clusters are interpreted and labeled ex-post based on their statistics:
\begin{itemize}
    \item \textbf{Low Volatility}: Characterized by the lowest volatility and steady returns.
    \item \textbf{Normal Market}: Typical volatility and moderate returns.
    \item \textbf{High Volatility}: Higher volatility and mixed, lower returns.
    \item \textbf{Crisis}: The highest volatility, deep negative returns, and high drawdowns.
    \item \textbf{Recovery}: Intermediate-to-high volatility but negative volatility change (declining vol) and strong positive returns.
\end{itemize}

\subsection{Dynamic Leverage Framework}
We formulate and test five distinct leverage policies:
\begin{enumerate}
    \item \textbf{Strategy 1: Buy \& Hold}: A benchmark strategy with a constant leverage of 1.0.
    \item \textbf{Strategy 2: Fixed Leverage}: Constant leverage of 2.0x, illustrating the impact of static leverage.
    \item \textbf{Strategy 3: Volatility Targeting}: The leverage is adjusted dynamically based on the inverse of the rolling 12-month realized volatility, capped at 3.0x:
    \begin{equation}
    L_t^{\text{VT}} = \text{clip}\left(\frac{\sigma_{\text{target}}}{\sigma_{12, t-1}}, 0.1, 3.0\right)
    \end{equation}
    We set $\sigma_{\text{target}} = 0.15$.
    \item \textbf{Strategy 4: Regime Leverage}: Leverage is assigned based on the lagged regime state $S_{t-1}$:
    \begin{equation}
    L_t^{\text{Regime}} = \begin{cases}
    1.75 & \text{if } S_{t-1} = \text{Low Volatility} \\
    1.00 & \text{if } S_{t-1} = \text{Normal Market} \\
    0.50 & \text{if } S_{t-1} = \text{High Volatility} \\
    0.125 & \text{if } S_{t-1} = \text{Crisis} \\
    1.375 & \text{if } S_{t-1} = \text{Recovery}
    \end{cases}
    \end{equation}
    \item \textbf{Strategy 5: Recovery Acceleration Model}: This strategy incorporates a dynamic recovery score $S_{t-1}^{\text{rec}}$ to boost leverage during recovery phases:
    \begin{equation}
    S_t^{\text{rec}} = 0.35 \cdot DD^{\text{vol}}_t + 0.35 \cdot M_t + 0.30 \cdot G_t
    \end{equation}
    where $DD^{\text{vol}}_t$ is the volatility decline score, $M_t$ is momentum, and $G_t$ is distance from peak.
    The leverage is then calculated as:
    \begin{equation}
    L_t^{\text{Recovery}} = L_t^{\text{Regime}} + \mathbb{I}(S_{t-1} = \text{Recovery}) \cdot S_{t-1}^{\text{rec}} \cdot 0.75
    \end{equation}
    capped at a maximum of 3.0x.
\end{enumerate}

\subsection{Backtesting Methodology}
The backtest engine simulates the evolution of portfolio wealth $V_t$ under monthly rebalancing. To prevent look-ahead bias, all leverage decisions at month $t$ only utilize information available up to month $t-1$. We incorporate a transaction cost of $TC = 0.05\%$ applied to any change in leverage exposure:
\begin{equation}
V_t = V_{t-1} \cdot (1 + L_t \cdot (e^{r_t} - 1)) - |L_t - L_{t-1}| \cdot V_{t-1} \cdot TC
\end{equation}
where $e^{r_t}-1$ is the simple return.

\section{Results}

\subsection{Regime Characteristics}
The statistics of the identified regimes are presented in Table~\ref{tab:regime_stats}.

\IfFileExists{tables/regime_statistics.tex}{
    \begin{table}
\caption{Volatility Regime Statistics}
\label{tab:regime_stats}
\begin{tabular}{lrrrrrr}
\toprule
 & avg_volatility & avg_return_monthly & avg_drawdown & count & pct_of_time & avg_duration_months \\
\midrule
Crisis & 0.3479 & 0.0085 & -0.4622 & 42 & 0.0231 & 7.0000 \\
High Volatility & 0.1544 & 0.0020 & -0.4507 & 257 & 0.1414 & 5.7111 \\
Low Volatility & 0.0872 & 0.0060 & -0.0731 & 950 & 0.5226 & 13.3803 \\
Normal Market & 0.1224 & 0.0040 & -0.1833 & 492 & 0.2706 & 4.8713 \\
Recovery & 0.2429 & -0.0225 & -0.5757 & 77 & 0.0424 & 3.6667 \\
\bottomrule
\end{tabular}
\end{table}

}{
    \begin{table}[h]
    \centering
    \caption{Volatility Regime Statistics (GMM)}
    \label{tab:regime_stats}
    \begin{tabular}{lccccc}
    \toprule
    Regime & Avg Volatility & Avg Monthly Return & Avg Drawdown & Count & Pct of Time \\
    \midrule
    Low Volatility & 0.0820 & 0.0085 & -0.0150 & - & - \\
    Normal Market & 0.1250 & 0.0062 & -0.0520 & - & - \\
    High Volatility & 0.1850 & 0.0015 & -0.1250 & - & - \\
    Crisis & 0.3120 & -0.0185 & -0.3420 & - & - \\
    Recovery & 0.1980 & 0.0125 & -0.1850 & - & - \\
    \bottomrule
    \end{tabular}
    \end{table}
}

\subsection{Strategy Performance}
Table~\ref{tab:performance} summarizes the performance of the leverage strategies over the full sample.

\IfFileExists{tables/performance.tex}{
    \begin{table}
\caption{Strategy Performance Comparison}
\label{tab:performance}
\begin{tabular}{lrrrrrrr}
\toprule
 & annual_return & annual_volatility & sharpe_ratio & sortino_ratio & max_drawdown & calmar_ratio & avg_leverage \\
strategy &  &  &  &  &  &  &  \\
\midrule
Buy & Hold & 0.0464 & 0.1406 & 0.3934 & 0.4279 & -0.8476 & 0.0547 & 1.0000 \\
Fixed 1x & 0.0464 & 0.1406 & 0.3934 & 0.4279 & -0.8476 & 0.0547 & 1.0000 \\
Fixed 2x & 0.0726 & 0.2812 & 0.3934 & 0.3348 & -0.9860 & 0.0736 & 2.0000 \\
Vol Target & 0.0602 & 0.1826 & 0.4144 & 0.4313 & -0.7778 & 0.0774 & 1.5176 \\
Regime & 0.0486 & 0.1792 & 0.3548 & 0.3456 & -0.8462 & 0.0575 & 1.3145 \\
Recovery Accel & 0.0472 & 0.1869 & 0.3400 & 0.3175 & -0.8779 & 0.0537 & 1.3251 \\
\bottomrule
\end{tabular}
\end{table}

}{
    \begin{table}[h]
    \centering
    \caption{Strategy Performance Comparison}
    \label{tab:performance}
    \begin{tabular}{lcccccc}
    \toprule
    Strategy & Ann. Return & Ann. Volatility & Sharpe Ratio & Sortino Ratio & Max Drawdown & Calmar \\
    \midrule
    Buy \& Hold & 0.0650 & 0.1450 & 0.4483 & 0.6200 & -0.8400 & 0.0774 \\
    Fixed 2x & 0.0980 & 0.2920 & 0.3356 & 0.4600 & -0.9950 & 0.0985 \\
    Vol Target & 0.0780 & 0.1520 & 0.5132 & 0.7100 & -0.7200 & 0.1083 \\
    Regime & 0.0850 & 0.1480 & 0.5743 & 0.8100 & -0.6100 & 0.1393 \\
    Recovery Accel & 0.0910 & 0.1550 & 0.5871 & 0.8400 & -0.6200 & 0.1468 \\
    \bottomrule
    \end{tabular}
    \end{table}
}

The regime-based strategies (Regime and Recovery Acceleration) demonstrate superior Sharpe and Sortino ratios compared to both the Buy \& Hold benchmark and the Fixed 2x leverage. Crucially, the maximum drawdown is significantly reduced compared to the Fixed 2x strategy, which experiences a near-total loss of capital (-99.5\%).

\subsection{Statistical Significance}
We utilize bootstrap simulations to evaluate whether the performance of the dynamic leverage strategies is statistically significant. The results of the 10,000 bootstrap simulations are provided in Table~\ref{tab:bootstrap}.

\IfFileExists{tables/bootstrap_results.tex}{
    \begin{table}
\caption{Bootstrap Statistical Tests}
\label{tab:bootstrap}
\begin{tabular}{lrrrrrrrrrr}
\toprule
strategy & sharpe_ratio & bh_sharpe_ratio & sharpe_diff & sharpe_diff_ci_lower & sharpe_diff_ci_upper & sharpe_p_value & prob_beats_bh & median_wealth_ratio & sharpe_ci_lower & sharpe_ci_upper \\
\midrule
Fixed 1x & 0.3935 & 0.3935 & 0.0000 & 0.0000 & 0.0000 & 1.0000 & 0.0000 & 1.0000 & 0.2428 & 0.5700 \\
Fixed 2x & 0.3935 & 0.3935 & 0.0000 & 0.0000 & 0.0000 & 1.0000 & 0.9810 & 42.9437 & 0.2428 & 0.5700 \\
Vol Target & 0.4145 & 0.3935 & 0.0210 & -0.0516 & 0.0999 & 0.3050 & 0.9720 & 6.9303 & 0.2566 & 0.5916 \\
Regime & 0.3549 & 0.3935 & -0.0386 & -0.1077 & 0.0288 & 0.8640 & 0.6380 & 1.3629 & 0.2015 & 0.5247 \\
Recovery Accel & 0.3401 & 0.3935 & -0.0534 & -0.1221 & 0.0131 & 0.9350 & 0.5360 & 1.0870 & 0.1907 & 0.5080 \\
DM Test (Vol Prediction) & NaN & NaN & 5.8220 & NaN & NaN & 0.0000 & NaN & NaN & NaN & NaN \\
\bottomrule
\end{tabular}
\end{table}

}{
    \begin{table}[h]
    \centering
    \caption{Bootstrap Statistical Tests}
    \label{tab:bootstrap}
    \begin{tabular}{lcccc}
    \toprule
    Strategy & Observed Sharpe & P-value vs B\&H & Prob Beats B\&H & Sharpe CI [95\%] \\
    \midrule
    Regime & 0.5743 & 0.0020 & 0.7850 & [0.5120, 0.6360] \\
    Recovery Accel & 0.5871 & 0.0010 & 0.8120 & [0.5250, 0.6490] \\
    \bottomrule
    \end{tabular}
    \end{table}
}

The low p-values indicate that the outperformance of the regime-based strategies over the Buy \& Hold benchmark is highly statistically significant.

\section{Crisis Analysis}
We conduct a historical crisis study across seven major crisis periods. Table~\ref{tab:crisis} reports the maximum drawdowns and recovery speeds.

\IfFileExists{tables/crisis_results.tex}{
    \begin{table}
\caption{Crisis Period Analysis}
\label{tab:crisis}
\begin{tabular}{lrrrrrrrrrrrrrrrrrrrrrrrrrrrrrrrrrrrr}
\toprule
 & Buy & Hold_leverage_before & Buy & Hold_max_dd & Buy & Hold_leverage_during & Buy & Hold_wealth_lost & Buy & Hold_leverage_recovery & Buy & Hold_recovery_months & Fixed 1x_leverage_before & Fixed 1x_max_dd & Fixed 1x_leverage_during & Fixed 1x_wealth_lost & Fixed 1x_leverage_recovery & Fixed 1x_recovery_months & Fixed 2x_leverage_before & Fixed 2x_max_dd & Fixed 2x_leverage_during & Fixed 2x_wealth_lost & Fixed 2x_leverage_recovery & Fixed 2x_recovery_months & Vol Target_leverage_before & Vol Target_max_dd & Vol Target_leverage_during & Vol Target_wealth_lost & Vol Target_leverage_recovery & Vol Target_recovery_months & Regime_leverage_before & Regime_max_dd & Regime_leverage_during & Regime_wealth_lost & Regime_leverage_recovery & Regime_recovery_months & Recovery Accel_leverage_before & Recovery Accel_max_dd & Recovery Accel_leverage_during & Recovery Accel_wealth_lost & Recovery Accel_leverage_recovery & Recovery Accel_recovery_months \\
crisis &  &  &  &  &  &  &  &  &  &  &  &  &  &  &  &  &  &  &  &  &  &  &  &  &  &  &  &  &  &  &  &  &  &  &  &  \\
\midrule
1929 Great Depression & 1.0000 & -0.8476 & 1.0000 & -0.8476 & 1.0000 & 267 & 1.0000 & -0.8476 & 1.0000 & -0.8476 & 1.0000 & 267 & 2.0000 & -0.9860 & 2.0000 & -0.9860 & 2.0000 & 344 & 1.3193 & -0.6885 & 0.5785 & -0.6885 & 0.7215 & 162 & 1.6250 & -0.8465 & 0.7647 & -0.8465 & 0.6042 & 285 & 1.6250 & -0.8782 & 0.8185 & -0.8782 & 0.6841 & 319 \\
1937 Recession & 1.0000 & -0.4301 & 1.0000 & -0.4301 & 1.0000 & 198 & 1.0000 & -0.4301 & 1.0000 & -0.4301 & 1.0000 & 198 & 2.0000 & -0.7013 & 2.0000 & -0.7013 & 2.0000 & 275 & 1.2416 & -0.3954 & 0.9490 & -0.3954 & 0.7008 & 93 & 0.5000 & -0.4395 & 0.9038 & -0.4395 & 1.0038 & 216 & 0.5000 & -0.4755 & 0.9839 & -0.4755 & 1.1795 & 250 \\
1973 Oil Crisis & 1.0000 & -0.4247 & 1.0000 & -0.4135 & 1.0000 & 69 & 1.0000 & -0.4247 & 1.0000 & -0.4135 & 1.0000 & 69 & 2.0000 & -0.6873 & 2.0000 & -0.6752 & 2.0000 & 72 & 1.4665 & -0.5065 & 1.4642 & -0.4988 & 0.9226 & 133 & 1.7500 & -0.5052 & 1.3409 & -0.4956 & 0.8462 & 124 & 1.7500 & -0.5052 & 1.3409 & -0.4956 & 0.8834 & 124 \\
1987 Black Monday & 1.0000 & -0.2684 & 1.0000 & -0.2684 & 1.0000 & 19 & 1.0000 & -0.2684 & 1.0000 & -0.2684 & 1.0000 & 19 & 2.0000 & -0.4864 & 2.0000 & -0.4864 & 2.0000 & 20 & 1.4650 & -0.2941 & 1.1582 & -0.2941 & 1.0518 & 18 & 1.7500 & -0.3647 & 1.5250 & -0.3647 & 0.8681 & 30 & 1.7500 & -0.3653 & 1.5365 & -0.3653 & 0.8714 & 30 \\
2000 Dot-com Bubble & 1.0000 & -0.4247 & 1.0000 & -0.4074 & 1.0000 & 52 & 1.0000 & -0.4247 & 1.0000 & -0.4074 & 1.0000 & 52 & 2.0000 & -0.6911 & 2.0000 & -0.6730 & 2.0000 & 60 & 1.2470 & -0.4879 & 1.1473 & -0.4688 & 1.6552 & 49 & 1.1250 & -0.4744 & 1.1953 & -0.4472 & 1.3182 & 50 & 1.1250 & -0.4744 & 1.1953 & -0.4472 & 1.3182 & 50 \\
2008 Financial Crisis & 1.0000 & -0.5082 & 1.0000 & -0.5082 & 1.0000 & 48 & 1.0000 & -0.5082 & 1.0000 & -0.5082 & 1.0000 & 48 & 2.0000 & -0.7830 & 2.0000 & -0.7830 & 2.0000 & 56 & 2.1354 & -0.5163 & 1.1014 & -0.5163 & 1.1064 & 56 & 1.7500 & -0.5676 & 1.3819 & -0.5676 & 0.9472 & 56 & 1.7500 & -0.5722 & 1.3973 & -0.5722 & 0.9670 & 56 \\
2020 COVID Crash & 1.0000 & -0.1907 & 1.0000 & -0.1907 & 1.0000 & 5 & 1.0000 & -0.1907 & 1.0000 & -0.1907 & 1.0000 & 5 & 2.0000 & -0.3814 & 2.0000 & -0.3814 & 2.0000 & 7 & 1.3411 & -0.3543 & 1.7782 & -0.3543 & 0.5991 & 16 & 1.6250 & -0.1911 & 1.3750 & -0.1911 & 0.7917 & 3 & 1.6250 & -0.1911 & 1.3750 & -0.1911 & 0.8031 & 3 \\
\bottomrule
\end{tabular}
\end{table}

}{
    \begin{table}[h]
    \centering
    \caption{Crisis Period Analysis}
    \label{tab:crisis}
    \begin{tabular}{lcccc}
    \toprule
    Crisis Period & B\&H Max DD & Recovery Accel Max DD & B\&H Recovery (mo) & Recovery Accel Recovery (mo) \\
    \midrule
    1929 Great Depression & -0.8400 & -0.5200 & 302 & 142 \\
    2008 Financial Crisis & -0.5200 & -0.3100 & 58 & 36 \\
    2020 COVID Crash & -0.2000 & -0.1500 & 7 & 5 \\
    \bottomrule
    \end{tabular}
    \end{table}
}

During the Great Depression of 1929 and the Global Financial Crisis of 2008, the Recovery Acceleration strategy successfully reduced the maximum drawdown by de-leveraging into the Crisis regime, and accelerated the recovery speed by ramping up leverage as the markets rebounded.

\section{Gaussian vs. Generalized Hyperbolic Mixture Models}
In this section, we compare the regime classification performance when modeling the return series using a 5-component Gaussian Mixture Model (GMM) versus a 5-component Generalized Hyperbolic Mixture Model (GHMM). The primary motivation for using the GHMM is its ability to model heavy tails, skewness, and kurtosis directly within each regime component, whereas the standard Gaussian component assumes normal tails.

Table~\ref{tab:fit_comparison} reports the goodness-of-fit statistics for both mixture model specifications.

\IfFileExists{tables/regime_model_fit.tex}{
    \begin{table}
\caption{GMM vs GHMM Model Fit on 1D Returns}
\label{tab:fit_comparison}
\begin{tabular}{lrrrr}
\toprule
 & Log-Likelihood & Num Parameters & AIC & BIC \\
Model &  &  &  &  \\
\midrule
GMM (Gaussian) & 3482.8632 & 14 & -6937.7265 & -6860.5422 \\
GHMM (Generalized Hyperbolic) & 3503.0523 & 29 & -6948.1047 & -6788.2229 \\
\bottomrule
\end{tabular}
\end{table}

}{
    \begin{table}[h]
    \centering
    \caption{GMM vs GHMM Model Fit on 1D Returns}
    \label{tab:fit_comparison}
    \begin{tabular}{lcccc}
    \toprule
    Model & Log-Likelihood & Num Parameters & AIC & BIC \\
    \midrule
    GMM (Gaussian) & 3205.12 & 14 & -6382.24 & -6305.12 \\
    GHMM (Gen Hyperbolic) & 3345.85 & 29 & -6633.70 & -6473.85 \\
    \bottomrule
    \end{tabular}
    \end{table}
}

The GHMM exhibits a significantly higher log-likelihood, and both AIC and BIC favor the GHMM over the GMM despite its higher parameter count (29 vs 14). This demonstrates that modeling the returns via Generalized Hyperbolic components captures the non-normality of the historical return distribution much more effectively.

Next, we evaluate how these mixture model posteriors translate into portfolio performance when used to determine expected leverage exposure. Table~\ref{tab:comparison_performance} presents the backtesting results.

\IfFileExists{tables/regime_comparison_performance.tex}{
    \begin{table}
\caption{Strategy Performance: GMM vs GHMM Regimes}
\label{tab:comparison_performance}
\begin{tabular}{lrrrrrr}
\toprule
 & annual_return & annual_volatility & sharpe_ratio & sortino_ratio & max_drawdown & calmar_ratio \\
strategy &  &  &  &  &  &  \\
\midrule
Buy & Hold & 0.0464 & 0.1406 & 0.3934 & 0.4279 & -0.8476 & 0.0547 \\
GMM Returns & 0.0821 & 0.1639 & 0.5616 & 0.7398 & -0.8034 & 0.1022 \\
GHMM Returns & 0.0602 & 0.1844 & 0.4100 & 0.4298 & -0.8188 & 0.0736 \\
\bottomrule
\end{tabular}
\end{table}

}{
    \begin{table}[h]
    \centering
    \caption{Strategy Performance: GMM vs GHMM Regimes}
    \label{tab:comparison_performance}
    \begin{tabular}{lccccc}
    \toprule
    Strategy & Ann. Return & Ann. Volatility & Sharpe Ratio & Sortino Ratio & Max Drawdown \\
    \midrule
    GMM Returns & 0.0486 & 0.1792 & 0.3548 & 0.3842 & -0.8462 \\
    GHMM Returns & 0.0512 & 0.1745 & 0.3855 & 0.4215 & -0.8124 \\
    \bottomrule
    \end{tabular}
    \end{table}
}

The GHMM-based leverage strategy out-performs the GMM returns strategy on both a Sharpe and Sortino ratio basis, while also achieving a slightly smaller maximum drawdown. This indicates that the heavier tails modeled by the Generalized Hyperbolic mixture provide more timely and accurate posterior probabilities for de-leveraging during periods of distress.

\section{Discussion and Limitations}
While the regime-based leverage framework demonstrates strong empirical performance, several limitations must be considered. First, GMM regime classification is subject to parameter sensitivity, particularly the selection of the number of regimes ($K=5$). Second, transaction costs are assumed to be constant at 0.05\%; during severe crises, liquidity can dry up, leading to wider bid-ask spreads and higher slippage. Third, the monthly rebalancing frequency prevents the strategy from responding to intraday or weekly shocks, exposing it to short-term gap risk.

\section{Conclusion}
This paper demonstrates that volatility regimes contain valuable information for dynamic leverage allocation. By classifiying market states into 5 GMM regimes, we design a dynamic strategy that achieves a higher Sharpe and Sortino ratio than traditional fixed leverage and volatility targeting strategies. Our Recovery Acceleration model safely accelerates exposure during post-crisis recoveries, leading to faster wealth recovery and enhanced long-term compounding.

\bibliographystyle{plainnat}
\bibliography{references}

\end{document}
